% Part: lambda-calculus
% Chapter: representability
% Section: currying

\documentclass[../../../include/open-logic-section]{subfiles}

\begin{document}

\olfileid{lam}{rep}{cur}
\olsection{کری‌سازی}

انتزاع~$\lambd$ای~$\lambd[x][M]$ تابعی تک‌آرگومانی را بازنمایی می‌کند؛
این امر هنگامی که می‌خواهیم تابعی چندآرگومانی تعریف کنیم محدودیتی جدی
است. یک راه آن بود که حساب~$\lambd$ را گسترش دهیم تا ساخت زوج‌ها،
سه‌تایی‌ها و مانند آن‌ها را مجاز کند؛ در آن صورت، برای نمونه، تابع
سه‌موضعی~$\lambd[x][M]$ انتظار داشت آرگومانش یک سه‌تایی باشد. اما انجام
این کار با \emph{کری‌سازی} مناسب‌تر است.

نمونه‌ای را در نظر بگیریم. لحظه‌ای وانمود می‌کنیم عمل~$+$ را در حساب
$\lambd$ در اختیار داریم. تابع جمع $2$-موضعی است، یعنی دو آرگومان
می‌گیرد. اما انتزاع~$\lambd$ تنها تابع تک‌آرگومانی به ما می‌دهد: نحو
عبارت‌هایی چون~$\lambd[(x,y)][(x+y)]$ را مجاز نمی‌کند. بااین‌حال، می‌توان
تابع تک‌موضعی~$f_x(y)$ را که با~$\lambd[y][(x+y)]$ داده می‌شود در نظر
گرفت؛ این تابع~$x$ را به آرگومان یگانهٔ خود~$y$ می‌افزاید. در واقع این
یک تابع منفرد نیست، بلکه خانواده‌ای از توابع متفاوت «افزودن~$x$» است،
یکی برای هر عدد~$x$. اکنون می‌توانیم تابع تک‌موضعی دیگری چون~$g$ را
به‌صورت~$\lambd[x][f_x]$ تعریف کنیم. با اعمال آن بر آرگومان~$x$،
$g(x)$ تابع~$f_x$ را برمی‌گرداند؛ پس ارزش‌هایش خود تابع‌اند. اکنون اگر
$g$ را بر~$x$ و سپس نتیجه را بر~$y$ اعمال کنیم، به این می‌رسیم:
$(g(x))y = f_x(y) = x+y$. بدین‌ترتیب، تابع تک‌موضعی~$g$ می‌تواند همان
کار تابع جمع دو‌موضعی را انجام دهد. «کری‌سازی» صرفاً نام این شگرد برای
تبدیل توابع دو‌موضعی به توابع تک‌موضعی است (که ارزش‌هایشان توابع
تک‌موضعی‌اند).

اینک نمونه‌ای را واقعاً در نحو حساب~$\lambd$ ببینیم. تابع
$f(x,y) = x$ را چگونه بازنمایی می‌کنیم؟ اگر بخواهیم تابعی تعریف کنیم که
دو آرگومان می‌گیرد و اولی را برمی‌گرداند، می‌توانیم
$\lambd[x][\lambd[y][x]]$ را بنویسیم؛ این عبارت به‌معنای دقیق کلمه تابعی
است که آرگومانی چون~$x$ می‌گیرد و تابع~$\lambd[y][x]$ را برمی‌گرداند.
تابع~$\lambd[y][x]$ آرگومان دیگری چون~$y$ می‌گیرد، اما آن را کنار
می‌گذارد و همیشه~$x$ را برمی‌گرداند. ببینیم وقتی
$\lambd[x][\lambd[y][x]]$ را بر دو آرگومان اعمال کنیم چه رخ می‌دهد:
\begin{align*}
  (\lambd[x][\lambd[y][x]])MN
  \bredone & (\lambd[y][M])N \\
  \bredone & M
\end{align*}

به‌طور کلی، برای نوشتن تابعی با پارامترهای~$x_1$، \dots،~$x_n$ که با
ترمی چون~$N$ تعریف شده است، می‌توانیم
$\lambd[x_1][\lambd[x_2][\ldots\lambd[x_n][N]]]$ را بنویسیم. اگر~$n$
آرگومان بر آن اعمال کنیم، به دست می‌آوریم:
\begin{multline*}
  (\lambd[x_1][\lambd[x_2][\ldots\lambd[x_n][N]]]) M_1 \dots M_n \bredone\\
  \begin{aligned}
  \bredone {} & (\Subst{(\lambd[x_2][\ldots\lambd[x_n][N]])}{M_1}{x_1}) M_2
  \dots M_n\\
   \eqs {} & (\lambd[x_2][\ldots\lambd[x_n][\Subst{N}{M_1}{x_1}]]) M_2
            \dots M_n \\
  \vdots & \\
  \bredone {} &\Subst{\Subst{N}{M_1}{x_1}\ldots}{M_n}{x_n}
  \end{aligned}
\end{multline*}
خط آخر به‌معنای دقیق کلمه یعنی~$M_i$ را به‌جای~$x_i$ در بدنهٔ تعریف
تابع جانشین کنیم؛ این دقیقاً همان چیزی است که هنگام اعمال چند آرگومان
بر یک تابع می‌خواهیم.
\end{document}

